\noindent{\sffamily\bfseries\color{brand}Abstract.}\hspace{0.4em}%
Vehicles are moving from a hundred function-specific electronic control
units (ECUs) to zonal electrical/electronic (E/E) architectures: a central
computer and a few zone controllers, linked by a Time-Sensitive Networking
(TSN) backbone, run most of the software. Before such a platform is built,
its designers must know how many cores each ECU needs, whether every
task-chain still meets its end-to-end deadline once its tasks are spread
over cores and over the network, and how much room is left for the
software updates of the coming years. Task allocation, fixed priorities,
shared-memory contention, Logical Execution Time (LET) communication
between cores and network delays are coupled, and, to our knowledge, no
existing method treats them together. This paper states the problem
formally and gives a system-level method that does. Task-chains are cut into sub-chains that
respect the bindings of sensors and actuators to ECUs, allocated to
heterogeneous cores by bin packing or by a new communication-aware rule,
consolidated onto fewer cores, and prioritized by an optimal assignment
that respects precedence. Response times include shared-memory
contention. The maximum reaction time of every task-chain is bounded,
under implicit communication or strict LET, with Network Calculus bounds for credit-based, time-aware and asynchronous
traffic shaping and for frame replication. Relocation and
deadline-driven moves, with an exact incremental re-analysis, repair
unschedulable cores and late task-chains. An exact constraint-programming
formulation with lazy schedulability cuts measures the optimality gap of the allocation,
and a lifecycle analysis adds software updates to the hardware fixed at
the start of production. On synthetic applications that follow published automotive
statistics, no simulated delay exceeds its bound, for 3377 network flows
and for 9000 task-chains simulated as a whole system. The heuristic is
2.3 to 2.9\,\% above the optimal allocation on average and proven optimal
in 81\,\% of 140 instances of up to 300 task-chains; on the default
platform it is as close as simulated annealing run for seconds, in
milliseconds, and closer than a genetic algorithm, which in turn
recovers most of the gap of the heuristic on harder platforms.
Communication awareness removes 64\,\% of the inter-ECU communications
for 5 more cores, strict LET shrinks the network bounds by 15--19\,\% but
costs 30 points of deadlines, and each core bought at the start of
production buys about half a software update. The same method maps domain, hybrid and
distributed architectures, within 1.1--1.2 and 4.4--4.7\,\% of the optimum on
the hybrid and domain platforms, and a whole-vehicle application of 2000
task-chains in about a second. \textbf{This project} originated from an idea conceived during the author's PhD
that could not be pursued at the time. It is an independent continuation
built with the help of AI tools, which assisted in developing the method,
code, and paper. The code and the scripts that reproduce every
experiment are at \url{https://github.com/ifupauline/mapping}.

\medskip\noindent{\sffamily\bfseries\color{brand}Keywords.}\hspace{0.4em}%
automotive embedded systems; software-defined vehicle; zonal E/E
architecture; real-time systems; task allocation; task-chains; cause-effect
chains;
logical execution time; fixed-priority scheduling; time-sensitive
networking; network calculus; asynchronous traffic shaping; frame
replication; memory contention; design-space exploration
